% ============================================================================
%  MCF-CAI State Vector Model
%  Theoretical Framework Document  ---  v1.0
%
%  Purpose: Formal definition of the MCF-CAI 3x3 object state grid and its
%           integration with IKK/MIR identity-information theory.
%
%  Status:  BLUE (newly introduced / provisional)
%  Compile: pdflatex (twice for ToC)
% ============================================================================
\documentclass[11pt, a4paper, oneside]{article}

\usepackage[top=1.0in, bottom=1.0in, left=1.1in, right=1.1in, headheight=14pt]{geometry}
\usepackage[T1]{fontenc}
\usepackage{lmodern}
\usepackage{microtype}
\usepackage{amsmath, amssymb, amsthm, mathtools}
\usepackage{bm}
\usepackage{mathrsfs}
\usepackage{booktabs, array, tabularx, longtable}
\usepackage{xcolor}
\usepackage{colortbl}
\usepackage{hyperref}
\usepackage{enumitem}
\usepackage{fancyhdr}
\usepackage{titlesec}

\hypersetup{colorlinks=true, linkcolor=black, urlcolor=blue, citecolor=black}

% ---- Colours ----------------------------------------------------------------
\definecolor{MacroBG}{HTML}{D4EDDA}
\definecolor{MacroFG}{HTML}{155724}
\definecolor{ChemBG}{HTML}{CCE5FF}
\definecolor{ChemFG}{HTML}{004085}
\definecolor{FundBG}{HTML}{F8D7DA}
\definecolor{FundFG}{HTML}{721C24}
\definecolor{CarrierBG}{HTML}{FFF3CD}
\definecolor{ActionBG}{HTML}{FFE0B2}
\definecolor{InfoBG}{HTML}{E2D9F3}

% ---- Theorem environments ---------------------------------------------------
\theoremstyle{definition}
\newtheorem{defn}{Definition}[section]
\newtheorem{axiom}{Axiom}[section]
\theoremstyle{plain}
\newtheorem{prop}{Proposition}[section]
\theoremstyle{remark}
\newtheorem{remark}{Remark}[section]
\newtheorem{warning}{Warning}[section]

% ---- MCF-CAI notation macros ------------------------------------------------

% Scale layers
\newcommand{\Macro}{\mathscr{M}}          % Macro layer
\newcommand{\Chem}{\mathscr{C}}           % Chemical layer
\newcommand{\Fund}{\mathscr{F}}           % Fundamental layer

% CAI basis
\newcommand{\Carrier}{_{\mathrm{C}}}      % Carrier/Configuration subscript
\newcommand{\Action}{_{\mathrm{A}}}       % Action/Coupling subscript
\newcommand{\Info}{_{\mathrm{I}}}         % Information/Entropy subscript

% Cell notation M_{s,b}
\newcommand{\MC}{\Macro_{\mathrm{C}}}
\newcommand{\MA}{\Macro_{\mathrm{A}}}
\newcommand{\MI}{\Macro_{\mathrm{I}}}
\newcommand{\CC}{\Chem_{\mathrm{C}}}
\newcommand{\CA}{\Chem_{\mathrm{A}}}
\newcommand{\CI}{\Chem_{\mathrm{I}}}
\newcommand{\FC}{\Fund_{\mathrm{C}}}
\newcommand{\FA}{\Fund_{\mathrm{A}}}
\newcommand{\FI}{\Fund_{\mathrm{I}}}

% Full state vector
\newcommand{\Vobj}{\mathcal{V}_{\mathrm{MCF}}^{\mathrm{CAI}}}

% IKK symbols (from Notation Registry)
\newcommand{\Dfrak}{\mathfrak{D}}
\newcommand{\Ifrak}{\mathfrak{I}}
\newcommand{\etaab}{\eta_{ab}}
\newcommand{\Psihid}{|\Psi^{\mathrm{hid}}\rangle}

% ---- Page style -------------------------------------------------------------
\pagestyle{fancy}
\fancyhf{}
\rhead{\small MCF-CAI State Vector Model v1.0}
\lhead{\small VSPER-SIM Theoretical Framework}
\cfoot{\thepage}

% ============================================================================
\begin{document}
% ============================================================================

\begin{center}
  {\LARGE\bfseries MCF-CAI State Vector Model}\\[0.5em]
  {\large Macro / Chemical / Fundamental $\;\times\;$ Carrier / Action / Information}\\[0.3em]
  {\normalsize Theoretical Framework Document v1.0 \quad|\quad VSPER-SIM v5.0.0-main}\\[0.5em]
  {\small Status: \colorbox{InfoBG}{BLUE --- Provisional / Newly Introduced}}
\end{center}

\vspace{0.5em}
\noindent\rule{\linewidth}{2pt}
\tableofcontents
\noindent\rule{\linewidth}{1pt}
\newpage

% ============================================================================
\section{Overview}
% ============================================================================

The MCF-CAI State Vector Model is a unified representational framework for
simulated objects in VSPER-SIM. It combines two orthogonal classification axes
into a single 3$\times$3 state grid:

\begin{itemize}
  \item \textbf{MCF} --- the vertical physical-layer axis: Macro, Chemical, Fundamental.
        This defines \emph{where} state is stored.
  \item \textbf{CAI} --- the horizontal identity-information basis: Carrier, Action, Information.
        This defines \emph{how} identity, coupling, and entropy-loss are tracked
        within each layer.
\end{itemize}

The compact equation form:
\[
  \boxed{
    \Vobj =
    \begin{bmatrix}
      \MC & \MA & \MI \\
      \CC & \CA & \CI \\
      \FC & \FA & \FI
    \end{bmatrix}
  }
\]

Or equivalently as a tensor sum:
\[
  \Vobj = \sum_{s \in \{\Macro,\Chem,\Fund\}} \;\sum_{b \in \{C,A,I\}} V_{s,b}\; \mathbf{e}_s \otimes \mathbf{e}_b
\]

where $\mathbf{e}_s$ are the scale-layer basis vectors and $\mathbf{e}_b$ are the
identity-basis vectors.

\begin{remark}
  MCF is the \emph{body plan}. CAI is the \emph{bookkeeping of identity and loss}.
  Neither axis can be reduced to the other. Both are required to fully specify
  the representational state of a simulated object.
\end{remark}

% ============================================================================
\section{The MCF Axis --- Physical Layer Hierarchy}
% ============================================================================

The MCF axis names three nested physical description layers. They are not disjoint:
a macro object contains chemical structure which contains fundamental degrees of
freedom. Each layer has its own observables, mediators, and resolution limit ---
i.e., each corresponds to a scale regime tuple $S_n = (D_n, M_n, L_n)$ from IKK IV.

\begin{longtable}{p{2cm} p{4cm} p{7cm}}
  \toprule
  \textbf{Layer} & \textbf{IKK scale} & \textbf{Contents} \\
  \midrule
  \rowcolor{MacroBG}
  $\Macro$ \quad Macro &
  $S_{\mathrm{mat}}$ (material scale) &
  Geometry, phase, grain structure, object boundary, bulk state,
  mechanical properties, continuum fields \\[4pt]
  \rowcolor{ChemBG}
  $\Chem$ \quad Chemical &
  $S_3$ (EM/spatial, atomistic) &
  Atoms, ions, molecules, ligands, bonds, coordination structure,
  element identity, orbital geometry \\[4pt]
  \rowcolor{FundBG}
  $\Fund$ \quad Fundamental &
  $S_0$--$S_2$ (existence, colour, subatomic) &
  Charge, spin proxy, nuclear state, isotope state, colour-force
  bookkeeping, hidden-coordinate content \\
  \bottomrule
\end{longtable}

\begin{warning}
  $\Fund$ does not imply full QCD. At the VSPER-SIM simulation scale,
  the fundamental layer is a \emph{proxy and bookkeeping layer} --- it tracks
  charge, spin intensity, and colour-averaged interaction (\texttt{cai\_channel})
  without a first-principles QCD kernel.
\end{warning}

% ============================================================================
\section{The CAI Axis --- Identity-Information Basis}
% ============================================================================

\begin{warning}
  \textbf{Namespace collision.}
  The abbreviation \texttt{CAI} already exists in VSPER-SIM source as
  \emph{Colour-Averaged Interaction} (force channel in \texttt{WO-VSEPR-SIM-EXTREME}).
  In theoretical documents, the CAI identity-information basis is always written
  in full or typeset in calligraphic/bold. The force channel is always written
  \texttt{cai\_channel} in monospace. These must never be conflated.
\end{warning}

The CAI axis defines three roles that any quantity can play within a given layer:

\begin{longtable}{p{2.5cm} p{3.5cm} p{7cm}}
  \toprule
  \textbf{Basis} & \textbf{Role} & \textbf{Contents} \\
  \midrule
  \rowcolor{CarrierBG}
  $C$ \quad Carrier &
  Configuration state &
  The \emph{what is here} component. Structural, configurational, geometric.
  This is the physical body of the object at this layer --- the vessel
  that holds all other quantities. \\[4pt]
  \rowcolor{ActionBG}
  $A$ \quad Action &
  Coupling / dynamics &
  The \emph{what is happening} component. Forces, reactions, flows,
  field couplings, phase transitions. The processes that change the
  carrier state over time. \\[4pt]
  \rowcolor{InfoBG}
  $I$ \quad Information &
  Entropy / identity loss &
  The \emph{what was lost} component. Formation history, projection loss,
  reaction pathway entropy, hidden-residual trace, recoverable identity
  fraction $\Dfrak$. This basis maps directly to the IKK sidecar layer. \\
  \bottomrule
\end{longtable}

The Information basis is the home of all IKK/MIR variables. It is never
used as a force input. It is a derived audit layer tracking what identity
content has been lost, preserved, or made unrecoverable as the simulation evolves.

% ============================================================================
\section{The Full 3x3 Grid --- Expanded Definitions}
% ============================================================================

\begin{longtable}{p{2.5cm} p{3.5cm} p{7cm}}
  \toprule
  \textbf{Cell} & \textbf{Name} & \textbf{Physical content} \\
  \midrule
  \rowcolor{MacroBG}
  $\MC$ & Macro Carrier &
  Position, orientation, phase, grain structure, object boundary,
  bulk geometry, bounding volume \\[3pt]
  \rowcolor{MacroBG}
  $\MA$ & Macro Action &
  Stress tensor, heat flux, flow velocity, fracture probability,
  diffusion rate, deformation gradient \\[3pt]
  \rowcolor{MacroBG}
  $\MI$ & Macro Information &
  Formation age, defect memory, measured macro state record,
  lost structural information, coarse-graining residual \\[6pt]
  \rowcolor{ChemBG}
  $\CC$ & Chemical Carrier &
  Atomic number $Z$, mass, bond topology, coordination structure,
  molecular geometry, ligand set \\[3pt]
  \rowcolor{ChemBG}
  $\CA$ & Chemical Action &
  Bond formation/breaking, reaction events, oxidation state change,
  solvation, ionisation, catalytic step \\[3pt]
  \rowcolor{ChemBG}
  $\CI$ & Chemical Information &
  Reaction pathway history, bond entropy loss, chemical distinguishability
  $\Dfrak_{\mathrm{chem}}$, formation route record \\[6pt]
  \rowcolor{FundBG}
  $\FC$ & Fundamental Carrier &
  Net charge $q$, spin proxy, nuclear state, isotope identity,
  colour-force bookkeeping vector \\[3pt]
  \rowcolor{FundBG}
  $\FA$ & Fundamental Action &
  Field coupling channels (EM, strong proxy, weak proxy), decay flag,
  annihilation event, projection event \\[3pt]
  \rowcolor{FundBG}
  $\FI$ & Fundamental Information &
  Hidden-coordinate residual $\Psihid$, projection loss $\Delta\Dfrak$,
  entropy-loss trace, unrecovered identity content \\
  \bottomrule
\end{longtable}

\begin{remark}
  $\FI$ is the deepest cell. It is the simulation-layer representation of
  $\mathfrak{I}_{\mathrm{true}} \setminus \mathfrak{I}_{\mathrm{top}}$ ---
  the identity content that exists at the primitive layer but cannot be
  recovered by any projection to the observable scale. In IKK IV terms:
  \[
    \FI \;\simeq\; |\Psi^{\mathrm{hid}}_{a\setminus b}\rangle
    = \bigl(\mathbf{I} - \Pi_{a\leftarrow b}\bigr)|\Psi_a\rangle
  \]
\end{remark}

% ============================================================================
\section{Relation to IKK Theory}
% ============================================================================

The CAI Information column is the simulation entry point for the full IKK
identity-information stack:

\begin{align*}
  \MI &\longleftrightarrow \text{IKK IV scale regime } S_{\mathrm{mat}}\text{: formation memory, coarse-graining loss} \\
  \CI &\longleftrightarrow \text{IKK II: chemical distinguishability } \Dfrak_{\mathrm{chem}}\text{, projection loss} \\
  \FI &\longleftrightarrow \text{IKK III: } \Psihid\text{, projection error } \varepsilon_{\mathfrak{I},ab}\text{, entropy trace}
\end{align*}

The IKK second law restated in MCF-CAI terms:

\[
  \Delta S > 0 \;\Longleftrightarrow\; \Delta\Dfrak < 0
  \;\Longleftrightarrow\;
  \text{Information column grows (more content lost)}
\]

The cross-scale projection chain maps directly onto the MCF sampling direction:

\[
  \Macro \xrightarrow{\;\text{resolve}\;} \Chem \xrightarrow{\;\text{resolve}\;} \Fund
  \;\longrightarrow\; \FI\text{ (projection loss accumulated)}
\]

And the reverse inference (defect propagation) maps the other direction:

\[
  \Fund\text{ perturbation} \xrightarrow{\;\text{propagate}\;} \Chem\text{ instability}
  \xrightarrow{\;\text{propagate}\;} \Macro\text{ defect}
\]

The $\Ifrak$-vector from WO-75B maps onto the CAI basis directly:

\begin{align*}
  I_{\mathrm{existence}} &\;\in\; \FC \\
  I_{\mathrm{EM}}        &\;\in\; \FC \;\cup\; \CA \\
  I_{\mathrm{spatial}}   &\;\in\; \CC \;\cup\; \MC \\
  I_{\mathrm{temporal}}  &\;\in\; \MA \;\cup\; \CA \\
  I_{\mathrm{internal}}  &\;\in\; \FC \\
  \Dfrak\;\text{(overall)}&\;\in\; \MI \;\cup\; \CI \;\cup\; \FI
\end{align*}

% ============================================================================
\section{VSIM Script Representation}
% ============================================================================

The VSIM object block directly encodes the 3$\times$3 grid as nested sections.
Each cell maps to a TOML-style sub-block:

\begin{verbatim}
[object]
name = "carbon_bead"

[object.macro.carrier]
position      = [0, 0, 0]
phase         = "solid"
geometry      = "sphere"

[object.macro.action]
stress        = [0, 0, 0]
heat_flux     = [0, 0, 0]

[object.macro.information]
formation_age = 0.0
defect_memory = 0.02

[object.chemical.carrier]
Z             = 6
mass_amu      = 12.011
bonds         = ["sigma", "pi"]

[object.chemical.action]
reaction_enabled  = true
oxidation_state   = 0
bond_exchange     = true

[object.chemical.information]
formation_route       = "graphitic_precursor"
reaction_entropy_loss = 0.0

[object.fundamental.carrier]
charge       = 0
spin_proxy   = 0.5
colour_state = [0.33, 0.33, 0.33]

[object.fundamental.action]
field_coupling = ["electric", "strong_proxy"]
decay_enabled  = false

[object.fundamental.information]
hidden_W        = 0.0
projection_loss = 0.0
entropy_loss    = 0.0
\end{verbatim}

\begin{remark}
  The \texttt{[object.*.information]} blocks are \textbf{sidecar-layer fields only}.
  They must never be used as force inputs or written back into ground-truth state.
  They are parsed into \texttt{IdentitySidecarRecord} derivatives, not into the
  particle kernel. This is the MCF-CAI encoding of the identity-sidecar doctrine
  from \texttt{identity\_sidecar.hpp}.
\end{remark}

% ============================================================================
\section{Notation Summary}
% ============================================================================

\begin{longtable}{p{3cm} p{4cm} p{6cm}}
  \toprule
  \textbf{Symbol} & \textbf{Name} & \textbf{Meaning} \\
  \midrule
  $\Vobj$   & MCF-CAI state vector  & Full 3$\times$3 object state grid \\
  $\Macro$  & Macro layer           & $S_{\mathrm{mat}}$ scale, bulk/continuum \\
  $\Chem$   & Chemical layer        & $S_3$ scale, atomistic/molecular \\
  $\Fund$   & Fundamental layer     & $S_0$--$S_2$, charge/spin/colour proxy \\
  $C$       & Carrier basis         & Configuration, structural state \\
  $A$       & Action basis          & Coupling, dynamics, process \\
  $I$       & Information basis     & Entropy, loss, identity audit \\
  $\MC$--$\FI$ & Grid cells (9)     & Each layer $\times$ each basis \\
  $\Dfrak$  & Distinguishability    & IKK Notation Registry; lives in $I$ column \\
  $\Psihid$ & Hidden residual       & Lives in $\FI$ \\
  $\Ifrak$  & Identity vector       & WO-75B; components distributed across $\FC$, $\CC$, $\MA$ \\
  \bottomrule
\end{longtable}

% ============================================================================
\section{Open Questions}
% ============================================================================

\begin{enumerate}
  \item \textbf{Off-diagonal couplings.}
    The 3$\times$3 grid treats each cell as independent. In reality,
    $\CI$ feeds back into $\MA$ (reaction entropy drives macro flow).
    A coupling tensor $T_{(s_1,b_1)(s_2,b_2)}$ is needed for full cross-cell dynamics.

  \item \textbf{Information column dynamics.}
    The $I$-column is currently static per frame. A time-evolution operator
    for the information column (analogous to $\hat{U}_a$ in IKK III for the
    $\Ifrak$-vector) is not yet defined.

  \item \textbf{Partial resolution.}
    A simulation may only have access to $\MC$ and $\CC$ (no fundamental
    bookkeeping). How does the grid degrade gracefully when lower rows
    are absent or null? A sparsity convention is needed.

  \item \textbf{CAI namespace.}
    The \texttt{CAI} abbreviation collision with \texttt{cai\_channel} (force)
    must be formally resolved in the Notation Registry. Proposed resolution:
    rename the force channel to \texttt{colour\_avg\_force} or \texttt{caf\_channel}
    in all source, freeing \texttt{CAI} for the identity basis exclusively.
\end{enumerate}

\end{document}